\documentclass{llncs}
\usepackage[utf8]{inputenc}
\usepackage[T1]{fontenc}
\usepackage[french]{babel}
\usepackage{graphicx}
\usepackage{amsmath,amssymb}
\usepackage{hyperref}
\usepackage{booktabs}

\begin{document}

\title{One-to-one SkewGRAM: de la biologie au NLP}
\author{
Bechir Mady\inst{1} \and 
Abderahmane Abderrahmane\inst{1} \and 
El Moustapha Mohamed El Moustapha\inst{1} \and 
Mohamedou Yahya Cheikh Mohamed Vall\inst{1}
}

\institute{
Institut Supérieur du Numérique (SupNum), Nouakchott, Mauritanie\\
\email{\{22014, 22053, 22013, 24239\}@supnum.mr}
}

\maketitle

\begin{abstract}
Le problème One-to-one SkewGRAM est un défi d'optimisation combinatoire complexe issu de la bio-informatique, visant à identifier des signatures métaboliques et génomiques conservées. Il nécessite de trouver un chemin dans un graphe orienté acyclique (représentant des voies métaboliques) tel que ses sommets induisent un sous-graphe connexe dans un graphe non-orienté (représentant des interactions génomiques). Les approches précédentes utilisant les réseaux de neurones sur graphes (GNN) ont montré des promesses mais peinent souvent à trouver des chemins longs ou nécessitent un post-traitement coûteux. Dans cet article, nous proposons une nouvelle approche inspirée du Traitement du Langage Naturel (NLP), spécifiquement le modèle Skip-Gram. Nous transposons cette technique aux graphes biologiques, en utilisant un processus de décodage en deux étapes et un échantillonnage négatif structurel. Notre méthode génère des chemins valides significativement plus rapidement que les solveurs exacts (ILP), atteignant une accélération moyenne de 68x, tout en offrant une alternative compétitive aux approches GNN existantes.
\keywords{Optimisation Combinatoire \and Traitement du Langage Naturel \and Skip-Gram \and Réseaux Biologiques \and Algorithmes de Graphes.}
\end{abstract}

\section{Introduction}
L'intelligence artificielle et ses applications sont de plus en plus pertinentes pour la bio-informatique, notamment dans les pays en développement où ces technologies offrent un potentiel de transformation rapide. Un défi majeur dans ce domaine est l'intégration de multiples réseaux biologiques pour extraire des modules fonctionnels significatifs. Le métabolisme cellulaire, par exemple, peut être modélisé comme un graphe orienté où les réactions biochimiques sont des arcs reliant des métabolites. Parallèlement, les interactions physiques entre les gènes ou les protéines peuvent être modélisées par des graphes non-orientés.

Le problème One-to-one SkewGRAM aborde l'intégration de ces deux niveaux d'information en cherchant des chemins qui sont cohérents à travers deux graphes différents : un graphe orienté (processus causaux) et un graphe non-orienté (associations fonctionnelles). L'identification de tels chemins permet de découvrir des modules fonctionnels conservés au cours de l'évolution.

Les avancées récentes ont appliqué les réseaux de neurones sur graphes (GNN) à ce problème, mais ces méthodes restent limitées soit par la courte longueur des chemins trouvés (en moyenne 1.2 nœuds), soit par la nécessité d'un post-traitement lourd. Dans ce travail, nous proposons d'améliorer l'approche GNN avec une méthode basée sur le NLP (Traitement du Langage Naturel) plus rapide. En adaptant le modèle Skip-Gram (popularisé par word2vec et DeepWalk), nous abordons ce problème spécifique de contraintes combinatoires.

\section{Formulation du Problème}
Le problème One-to-one SkewGRAM peut être défini formellement comme suit :

Soit $V = \{v_1, v_2, \dots, v_n\}$ un ensemble de nœuds (entités biologiques).
\begin{itemize}
    \item $D = (V, A)$ est un graphe orienté acyclique (DAG) représentant des processus biologiques directionnels.
    \item $G = (V, E)$ est un graphe non-orienté encodant des interactions fonctionnelles.
\end{itemize}

L'objectif est de trouver un chemin $P = (v_1, \dots, v_k)$ dans $D$ tel que :
\begin{enumerate}
    \item \textbf{Condition de chemin} : $(v_i, v_{i+1}) \in A$ pour tout $i$, c'est-à-dire que $P$ est un chemin orienté valide dans $D$.
    \item \textbf{Condition de connectivité} : Le sous-graphe de $G$ induit par les sommets $\{v_1, \dots, v_k\}$ est connexe.
\end{enumerate}

Un tel chemin est dit \textbf{(D,G)-consistant}. Trouver le plus long chemin (D,G)-consistant est un problème NP-difficile. Les solveurs exacts comme la Programmation Linéaire en Nombres Entiers (ILP) fournissent la longueur optimale mais souffrent de temps d'exécution exponentiels sur de grands graphes.

\section{Approche NLP Proposée : Skew-GRAM}
Notre approche s'inspire des succès du NLP pour l'apprentissage de représentations, en traduisant les nœuds de graphe en "mots" et les marches aléatoires sur $D$ en "phrases".

\subsection{Plongements de Nœuds (Embeddings) et Skip-Gram}
À l'instar de DeepWalk ou node2vec, nous générons des séquences de nœuds via des marches aléatoires sur le graphe orienté $D$. Nous apprenons ensuite des matrices d'embeddings $\mathbf{V}_{in}, \mathbf{V}_{out} \in \mathbb{R}^{N\times d}$ ($d=32$). 
Cependant, contrairement au Skip-Gram classique, nous introduisons un \textbf{échantillonnage négatif structurel}.

La fonction de perte est optimisée comme suit :
\begin{equation}
\mathcal{L}(c,o,n_{1:K}) = -\log\sigma(\mathbf{v}_{in,c}^\top \mathbf{v}_{out,o}) - \sum_{k=1}^K \log\sigma(-\mathbf{v}_{in,c}^\top \mathbf{v}_{out,n_k})
\end{equation}

Les négatifs structurels $n_k$ sont tirés parmi les non-successeurs du nœud cible $c$ dans $D$ ($P_{struct}(n\mid c) \propto \mathbf{1}[(c,n)\notin A]$). Cela force l'espace latent à pénaliser fortement les transitions impossibles dans le graphe causal, capturant ainsi rigoureusement la structure du DAG.

\begin{figure}[htbp]
\centering
\includegraphics[width=0.8\textwidth]{figures/04_pca_embeddings.png}
\caption{Projection PCA des embeddings de nœuds générés par Skew-GRAM. La structure latente capture l'ordre topologique du graphe orienté $D$.}
\label{fig:pca}
\end{figure}

Nous appliquons également un sous-échantillonnage des "hubs" (nœuds très connectés) pour éviter que les chemins très denses ne dominent les probabilités de transition.

\subsection{Décodage en Deux Étapes}
Étant donné que la condition de connectivité dans $G$ s'applique à l'ensemble \emph{final} des sommets du chemin (qui peut induire un arbre, et non forcément un chemin séquentiel dans $G$), il est impossible de garantir la connexité de manière incrémentale.
Nous utilisons donc un décodeur en deux étapes :

\begin{enumerate}
    \item \textbf{Étape A (Génération)} : Recherche en faisceau (beam search, largeur 4) sur les successeurs de $D$, en masquant les nœuds déjà visités. La probabilité de sélection est guidée par la similarité des embeddings appris :
    $$P_{skew}(v\mid u,\mathcal{U}) = \frac{\mathbf{1}[(u,v)\in A]\cdot\mathbf{1}[v\notin\mathcal{U}]\cdot \exp(\mathrm{sim}(u,v)/\tau)}{\sum_{v'\in N^+(u)\setminus\mathcal{U}}\exp(\mathrm{sim}(u,v')/\tau)}$$
    \item \textbf{Étape B (Filtrage de Connexité)} : Sur le chemin candidat (de longueur $L$), nous appliquons une fenêtre glissante. Pour chaque sous-intervalle contigu, nous vérifions si l'ensemble des sommets induit un graphe connexe dans $G$ en utilisant un algorithme Union-Find incrémental en $O(L^2\cdot\alpha(L))$. Nous extrayons ainsi le plus long sous-chemin $(D,G)$-consistant.
\end{enumerate}

\begin{figure}[htbp]
\centering
\includegraphics[width=0.9\textwidth]{figures/06_graph_topology_paths.png}
\caption{Visualisation de la topologie. À gauche : le chemin candidat généré dans le graphe $D$. À droite : le sous-graphe induit dans le graphe $G$, illustrant la connectivité finale.}
\label{fig:topology}
\end{figure}

\section{Résultats Expérimentaux}
Nous avons évalué notre approche sur un ensemble d'instances synthétiques d'Erdős-Rényi (environ 100 nœuds par instance), qui modélisent les caractéristiques des réseaux biologiques. 
Nous avons comparé les longueurs de chemin et les temps d'exécution de notre méthode Skew-GRAM avec le solveur exact ILP2.

\begin{figure}[htbp]
\centering
\includegraphics[width=0.9\textwidth]{figures/03_evaluation_comparison.png}
\caption{À gauche : Comparaison des longueurs de chemins par instance (ILP2 vs Skew-GRAM). À droite : Compromis qualité/temps, démontrant la grande rapidité de Skew-GRAM.}
\label{fig:results}
\end{figure}

Les résultats démontrent que la méthode Skew-GRAM est \textbf{extrêmement rapide}. Alors que ILP2 prend en moyenne 28 secondes par instance, Skew-GRAM génère le chemin en 0.41 seconde, soit une \textbf{accélération de plus de 68x}.
En termes de qualité, le chemin généré lors de l'étape A possède une longueur moyenne de 6.45. Après le filtrage sévère de l'étape B pour garantir la $(D,G)$-consistance, la longueur moyenne est de 3.42 (contre 9.97 pour la solution optimale ILP2).
Bien que la méthode n'atteigne pas encore la longueur optimale exacte, elle génère des chemins valides (contrairement aux méthodes de Machine Learning traditionnelles comme Random Forest) et de longueur très supérieure au 1.2 nœuds obtenus par des décodeurs GNN séquentiels stricts.

\section{Conclusion et Travaux Futurs}
En traitant les réseaux métaboliques et génomiques comme des séquences linguistiques et en exploitant le modèle Skip-Gram avec échantillonnage négatif structurel, nous avons introduit une heuristique très rapide pour le problème NP-difficile One-to-one SkewGRAM. Cette méthode se positionne comme un excellent filtre de pré-traitement (warm-start) pour les méthodes exactes, surclassant les GNN traditionnels en termes de vitesse d'inférence.

Les travaux futurs viseront à intégrer directement la connectivité du graphe $G$ lors de la génération (Étape A) via un entraînement conjoint (Joint Embeddings), afin de réduire l'écart de longueur de chemin avec les solveurs exacts.

\section*{Remerciements}
Nous remercions chaleureusement nos superviseurs, le Dr. Mohamed Lemine Ahmed Sidi et le Dr. Hafedh Mohamed Babou, pour leurs orientations scientifiques tout au long de ce projet.

\begin{thebibliography}{10}

\bibitem{mikolov}
Mikolov, T., Sutskever, I., Chen, K., Corrado, G., Dean, J.:
Distributed Representations of Words and Phrases and their Compositionality.
NeurIPS (2013)

\bibitem{deepwalk}
Perozzi, B., Al-Rfou, R., Skiena, S.:
DeepWalk: Online Learning of Social Representations.
KDD (2014)

\bibitem{node2vec}
Grover, A., Leskovec, J.:
node2vec: Scalable Feature Learning for Networks.
KDD (2016)

\bibitem{sidi}
Ahmed Sidi, M. L., Babou, H. M., Dhib, C., et al.:
Constraint programming approaches for finding conserved metabolic and genomic patterns.
Computers \& Operations Research (2025)

\end{thebibliography}

\end{document}
